\documentclass[11pt,reqno]{amsart}
\usepackage[T1]{fontenc}
\usepackage{iftex}
\ifPDFTeX
  \usepackage[utf8]{inputenc}
\fi
\usepackage{lmodern}
\usepackage{amsmath,amssymb,amsthm,mathtools,mathrsfs}
\usepackage{booktabs}
\usepackage{enumitem}
\usepackage{microtype}
\usepackage{xcolor}
\usepackage[hidelinks]{hyperref}
\usepackage[nameinlink,noabbrev]{cleveref}

\numberwithin{equation}{section}
\raggedbottom

\newtheorem{theorem}{Theorem}[section]
\newtheorem{proposition}[theorem]{Proposition}
\newtheorem{lemma}[theorem]{Lemma}
\newtheorem{corollary}[theorem]{Corollary}
\newtheorem{question}[theorem]{Question}
\theoremstyle{definition}
\newtheorem{definition}[theorem]{Definition}
\newtheorem{example}[theorem]{Example}
\theoremstyle{remark}
\newtheorem{remark}[theorem]{Remark}

\newcommand{\A}{\mathbb A}
\newcommand{\C}{\mathbb C}
\newcommand{\F}{\mathbb F}
\newcommand{\Gm}{\mathbb G_{\mathrm m}}
\newcommand{\Lef}{\mathbb L}
\newcommand{\PP}{\mathbb P}
\newcommand{\Cl}{\operatorname{Cl}}
\newcommand{\Disc}{\operatorname{Disc}}
\newcommand{\Res}{\operatorname{Res}}
\newcommand{\Sing}{\operatorname{Sing}}
\newcommand{\Sym}{\operatorname{Sym}}
\newcommand{\Spec}{\operatorname{Spec}}
\newcommand{\ord}{\operatorname{ord}}
\newcommand{\id}{\operatorname{id}}
\newcommand{\cE}{E_{\mathrm c}}
\newcommand{\cchi}{\chi_{\mathrm c}}
\newcommand{\set}[1]{\left\{#1\right\}}
\newcommand{\abs}[1]{\left|#1\right|}
\newcommand{\angles}[1]{\left\langle#1\right\rangle}

\crefname{theorem}{theorem}{theorems}
\Crefname{theorem}{Theorem}{Theorems}
\crefname{proposition}{proposition}{propositions}
\Crefname{proposition}{Proposition}{Propositions}
\crefname{lemma}{lemma}{lemmas}
\Crefname{lemma}{Lemma}{Lemmas}
\crefname{corollary}{corollary}{corollaries}
\Crefname{corollary}{Corollary}{Corollaries}
\crefname{question}{question}{questions}
\Crefname{question}{Question}{Questions}

\usepackage{mathrsfs}

\title[Boundary Belyi covers and a degree-296 obstruction]{Boundary Belyi Covers and a Degree-296
Normal-Jet Obstruction in the Plane Jacobian Problem}
\author{Nathaniel Monson}
\date{July 24, 2026, revised proof draft}

\begin{document}

\begin{abstract}
Newton-face differential equations are a classical tool in the plane
Jacobian problem.  We isolate the primitive three-point covers carried by a
terminal face and the filtered obstruction theory carried by its formal
normal neighborhood.

A terminal equation
\[
 npq-mzpq'+nzp'q=1
\]
defines a Belyi map.  When the polynomial exponent lattice has gap \(g\),
the relevant map is not the ambient degree-\(mnb\) cover but its
\(\mu_g\)-quotient of degree \(mnb/g\).  We give its passport, prove
automatic connectedness, and compute the first five quotient Hurwitz
problems in the complete-chain queue: they have \(1,1,1,2,2\) classes.
The explicit quotient points through maximum degree \(128\) are
infinitesimally rigid modulo source scaling.

At every normal order the new coefficients enter through one universal
first-order operator.  We put it in logarithmic gauge normal form, identify
the dual of each finite Newton-window cokernel with filtered principal
parts, and reconstruct the exact compatibility equations as residue
pairings.  The normalized adjacent-contact equation recovers a classical
truncated-binomial Shabat map, now forced by the boundary-volume transport.

For the exact terminal full-support \((8,28)\) model, five selected
filtered-residue equations define a finite \'{e}tale scheme of degree \(296\)
over the quintic lower-face field, and the sixth selected obstruction is a
unit on that scheme.  The proof combines an exact mixed-volume calculation,
a complete audit of all \(344\) proper toric faces, a \(296\)-dimensional
finite-algebra certificate, and characteristic-zero lifting from a good
split prime.  Thus the six exact obstruction polynomials have no common
zero.  Interpreting this as a theorem about every terminal Keller model
still depends on the separately stated completeness and normalization
inputs; no degree-\(125\) theorem or proof of the plane Jacobian conjecture
is claimed.
\end{abstract}

\maketitle

\begin{center}
\small\emph{Working computer-assisted proof draft.  The recovered exact
certificates have been rerun.  Independent implementation and specialist
review remain submission requirements.}
\end{center}

\section{Introduction}
\label{sec:introduction}

The counterexample to the Jacobian conjecture in dimension three leaves the
two-variable conjecture open.  The separation is geometric: the marked-root
construction that fills an affine threefold does not directly fill an
affine plane.  In the plane, the unresolved information accumulates at
infinity in Newton faces, divisorial valuations, and compatibility between
successive boundary layers.

We credit Akhil Mathew as the primary human source of the problem: Mathew
prompted Levent Alp\"oge, Alp\"oge prompted Fable, Fable produced the
three-dimensional example, and Alp\"oge announced it publicly
\cite{ulam2026counterexample,alpoge2026announcement}.

The background here is substantial.  One-variable face differential
equations of the form used below occur in the work of Formanek and Stothers
on the plane Jacobian problem \cite{formanek2011stothers}.  Borisov developed
Belyi maps and dessins in compactification frameworks for hypothetical
plane counterexamples \cite{borisov2009geometric,borisov2020frameworks}.  Guccione, Guccione,
Horruitiner, and Valqui constructed complete-chain algorithms and reduced
the range below maximum degree \(125\) to the pair \((72,108)\) and its
transpose \cite{guccioneEtAl2017algorithms,guccioneEtAl2022degree108}.
Normal-boundary determinant expansions also have a recent valuative
precursor in Bado's work \cite{bado2026valuative}.

While this manuscript was being prepared, a MathOverflow answer announced a
computer-assisted calculation eliminating the remaining supports below
maximum coordinate degree \(125\), with a write-up in preparation
\cite{ratto3423degree125}.  A separate
exact-arithmetic campaign by Santibáñez-Leal states strong certificate
evidence at the same frontier while distinguishing that evidence from a
complete all-coefficient proof \cite{santibanezLeal2026planar}.  We do not
claim the degree-\(125\) bound.  The purpose of this paper is to extract
reusable boundary theorems from an independent terminal calculation and to
state exactly what remains missing from a global argument.

\subsection{Main results}

The first result corrects a tempting overcount.  A terminal face may be
written in a fractional uniformizing variable \(z\), but the exponent
lattice of the original polynomial pair can force both face polynomials to
belong to \(k[z^g]\).  The actual Hurwitz problem is then a quotient by
\(u=z^g\).

\begin{theorem}[Lattice-gap quotient, informal form]
\label{thm:intro-primary}
Let \(m,n>1\) be coprime, let \(b\ge1\), and suppose a terminal face satisfies
\[
npq-mzpq'+nzp'q=1,\qquad
\deg p=mb-1,\quad \deg q=nb.
\]
If \(p(z)=\bar p(z^g)\) and \(q(z)=\bar q(z^g)\), as forced by a
complete-chain lattice gap \(g\), then \(g\mid n\) and
\[
\bar\tau(u)=u^{n/g}\frac{\bar p(u)^n}{\bar q(u)^m}
\]
is a Belyi map of degree \(mnb/g\), with passport
\[
\left(n^{(mb-1)/g},\frac ng\right),\qquad
\left(m^{nb/g}\right),\qquad
\left(\frac{(m+n)b-1}{g},
1^{\frac{mnb-(m+n)b+1}{g}}\right).
\]
Every permutation triple with this passport is transitive.
\end{theorem}

Applying the theorem to the first five terminal complete-chain rows gives
quotient degrees \(6,10,9,9,16\) and respectively \(1,1,1,2,2\) dessin
classes.  The first row is especially useful: an ambient degree-\(30\)
passport has eleven classes, but the gap-five lattice selects the unique
cyclic pullback of
\[
\frac{u(u-1)^5}{(u^2-\frac53u+\frac59)^3}.
\]

The second result describes the normal direction.  At layer \(r\), new
coefficients enter the determinant equation through a universal first-order
operator
\[
\mathscr D_r^{\alpha,\beta}(a,b)
=(\alpha-r)a\,dB_0-\beta B_0\,da
+\alpha A_0\,db+(r-\beta)b\,dA_0.
\]
After passing to relative variations, this becomes a rank-one logarithmic
connection plus multiplication by the fixed boundary-volume form.  For the
exact two-point Newton windows, every coefficient left-null vector is
canonically a filtered principal part, and every compatibility polynomial
is its residue pairing with the lower-layer forcing term.

The adjacent reduced cover forced by normalized contact is the following
classical Shabat family, recovered here from boundary transport.

\begin{theorem}[Secondary transport, informal form]
\label{thm:intro-secondary}
For integers \(a\ge1\) and \(e>a^2\), set \(\delta=e-a^2\).  The normalized
boundary-volume transport equation has the unique polynomial-numerator
solution
\[
h(s)=\frac{H_{a,e}(s)}{(s-1)^a},
\qquad
H_{a,e}(s)=
\sum_{k=0}^a
\frac{a^k a!}{(a-k)!\prod_{j=0}^k(e-aj)}s^k.
\]
The exceptional ratio
\[
W_{a,e}(s)=\frac{(s-1)^\delta H_{a,e}(s)^a}{s^e}
\]
is a degree-\(e\) Belyi map with passport
\[
(a^a,\delta),\qquad(e),\qquad(a+1,1^{e-a-1}).
\]
\end{theorem}

After the reciprocal change \(s=1/x\), this is the truncated-binomial
Shabat polynomial appearing in the hypergeometric Belyi literature
\cite{vidunas2025belyi}.  The new point here is that boundary-volume
transport forces this classical map and identifies its parameters with
adjacent divisorial data.  For \((a,e)=(4,17)\), it recovers the degree-
\(17\) secondary cover in the terminal \((8,28)\) calculation.

The final result is a finite no-gluing theorem for the exact terminal model.

\begin{theorem}[Degree-\(296\) terminal obstruction, informal form]
\label{thm:intro-degree-296}
Over the quintic lower-face field, five selected filtered-residue equations
define a finite \'{e}tale scheme \(Z\) of degree \(296\) in the five-parameter
torus.  The remaining selected residue coordinate \(\rho\) is a unit on
\(Z\).  Equivalently,
\[
 \operatorname{Nm}_{Z/K}(\rho)\ne0,
\]
and the six exact obstruction polynomials have no common zero.
\end{theorem}

The proof audits every proper face of the Minkowski sum of the five Newton
polytopes, establishes BKK nondegeneracy at a good split prime, identifies
the complete \(296\)-dimensional reduced special algebra, and lifts the
finite \'{e}tale scheme and the norm obstruction to characteristic zero.  The
geometric lesson is that the reduced boundary covers and several initial
normal compatibilities coexist: the contradiction appears only in the
global Newton-bounded thickening.

\subsection{Scope}

The lattice-gap, Hurwitz, filtered-principal-part, mixed-volume, toric-face,
and degree-\(296\) statements are theorems about the explicitly defined
face and normal-jet systems.  In particular, the six exact terminal
obstruction polynomials have no common zero over an algebraic closure of the characteristic-zero field.  We do not
prove:

\begin{itemize}
\item that the imported complete-chain and support reductions exhaust every
hypothetical Keller pair below degree \(125\);
\item that the stated normalizations and saturations cover every terminal
vertex-saturated case;
\item a general all-layer algebraization or approximate-root descent theorem;
or
\item the plane Jacobian conjecture.
\end{itemize}

These dependency boundaries remain visible throughout the paper.

\section{Face equations and their Belyi maps}
\label{sec:faces}

We work over a characteristic-zero field \(k\), embedded in \(\C\) when
using topological language about covers.  A passport of a degree-\(D\)
Belyi map is the ordered triple of partitions of \(D\) recording
ramification over the three branch values.  The order of the branch values
is immaterial for our applications.

\subsection{The ambient terminal face}

Let \(m,n>1\) be coprime and \(b\ge1\).  Consider polynomials
\[
p,q\in k[z],\qquad
\deg p=mb-1,\quad \deg q=nb,
\]
satisfying
\begin{equation}
\label{eq:ambient-ode}
npq-mzpq'+nzp'q=1.
\end{equation}
Define
\begin{equation}
\label{eq:ambient-tau}
\tau(z)=z^n\frac{p(z)^n}{q(z)^m}.
\end{equation}

\begin{proposition}[Face-to-passport formula]
\label{prop:ambient-passport}
The rational function \(\tau\) is a degree-\(mnb\) Belyi map and
\begin{equation}
\label{eq:ambient-derivative}
\tau'(z)=z^{n-1}\frac{p(z)^{n-1}}{q(z)^{m+1}}.
\end{equation}
Its passport is
\begin{equation}
\label{eq:ambient-passport}
\left(n^{mb}\right),\qquad
\left(m^{nb}\right),\qquad
\left((m+n)b-1,1^{mnb-(m+n)b+1}\right).
\end{equation}
\end{proposition}

\begin{proof}
Logarithmic differentiation of \eqref{eq:ambient-tau} gives
\[
\tau'
=z^{n-1}\frac{p^{n-1}}{q^{m+1}}
\bigl(npq+nzp'q-mzpq'\bigr).
\]
Equation \eqref{eq:ambient-ode} yields
\eqref{eq:ambient-derivative}.

At \(z=0\), the face equation gives \(np(0)q(0)=1\).  At a zero
\(z_0\) of \(p\), it gives \(nz_0p'(z_0)q(z_0)=1\); hence the root is
simple, nonzero, and not a root of \(q\).  The analogous assertion follows
for roots of \(q\).  Thus the zero fiber of \(\tau\) has one point of
multiplicity \(n\) at zero and \(mb-1\) further points of multiplicity
\(n\).  The pole fiber consists of \(nb\) points of multiplicity \(m\).

The numerator and denominator of \(\tau\) are coprime and both have degree
\(mnb\), so \(\tau(\infty)\in k^\times\).  From
\eqref{eq:ambient-derivative},
\[
\tau'(z)=O\!\left(z^{-(m+n)b}\right)
\quad\text{as }z\longrightarrow\infty.
\]
Therefore
\[
\tau(z)-\tau(\infty)
=O\!\left(z^{-((m+n)b-1)}\right),
\]
with nonzero leading coefficient.  The ramification contributions already
identified total
\begin{align*}
mb(n-1)+nb(m-1)+((m+n)b-2)
=2mnb-2.
\end{align*}
Riemann--Hurwitz leaves no additional critical point.  The rest of the
\(\tau(\infty)\)-fiber is unramified, proving
\eqref{eq:ambient-passport}.
\end{proof}

\begin{remark}[Relation with Formanek--Stothers]
Putting \(P=zp\) turns \eqref{eq:ambient-ode} into
\[
 nP'q-mPq'=1.
\]
Thus the ambient equation is exactly in the classical Formanek--Stothers
extra-special-pair family \cite{formanek2011stothers}.  The new input below
is the primitive lattice quotient, its forced passport and connectedness,
and the exact terminal Hurwitz calculations.
\end{remark}

\subsection{The lattice quotient}

In a complete-chain chart, \(z\) can be a fractional uniformizing
coordinate.  The polynomial exponent lattice forces a gap
\[
g=\frac{\rho}{\gcd(\rho,\ell)}
\]
and restricts the face polynomials to \(k[z^g]\).  The following theorem is
stated purely algebraically; the complete-chain theory supplies its
hypothesis in the applications.

\begin{theorem}[Lattice-gap terminal cover]
\label{thm:lattice-gap}
Assume the hypotheses of \cref{prop:ambient-passport} and suppose
\[
p(z)=\bar p(u),\qquad q(z)=\bar q(u),\qquad u=z^g,
\]
for an integer \(g\ge1\).  Then
\[
g\mid mb-1,\qquad g\mid nb,\qquad \gcd(g,b)=1,\qquad g\mid n.
\]
Put
\begin{align*}
N&=\frac ng,&
A&=\frac{mb-1}{g},&
B&=\frac{nb}{g},\\
D&=\frac{mnb}{g},&
H&=\frac{(m+n)b-1}{g}.&&
\end{align*}
The quotient face equation is
\begin{equation}
\label{eq:quotient-ode}
N\bar p\bar q-mu\bar p\bar q'
+nu\bar p'\bar q=\frac1g,
\end{equation}
and
\begin{equation}
\label{eq:quotient-map}
\bar\tau(u)=u^N\frac{\bar p(u)^n}{\bar q(u)^m}
\end{equation}
is a degree-\(D\) Belyi map with
\begin{equation}
\label{eq:quotient-derivative}
\bar\tau'(u)=\frac1g
u^{N-1}\frac{\bar p(u)^{n-1}}{\bar q(u)^{m+1}}
\end{equation}
and passport
\begin{equation}
\label{eq:quotient-passport}
\left(n^A,N\right),\qquad
\left(m^B\right),\qquad
\left(H,1^{D-H}\right).
\end{equation}
Moreover,
\[
\tau(z)=\bar\tau(z^g),
\]
and every permutation triple with passport
\eqref{eq:quotient-passport} is transitive.
\end{theorem}

\begin{proof}
The degree assumptions imply \(g\mid mb-1\) and \(g\mid nb\).
The first divisibility gives \(\gcd(g,b)=1\), so the second gives \(g\mid n\).
Substitution into \eqref{eq:ambient-ode} and division by \(g\) yield
\eqref{eq:quotient-ode}; logarithmic differentiation gives
\eqref{eq:quotient-derivative}.  The proof of
\cref{prop:ambient-passport}, with the zero multiplicity \(N\) at \(u=0\),
gives the degree and passport.  The factorization of \(\tau\) is immediate.

It remains to prove transitivity.  Let
\((\sigma_0,\sigma_1,\sigma_\infty)\) have the stated cycle types and
product one.  An orbit disjoint from the unique \(H\)-cycle of
\(\sigma_\infty\) is pointwise fixed by \(\sigma_\infty\).  On that orbit,
\(\sigma_1=\sigma_0^{-1}\), so its \(\sigma_0\)- and
\(\sigma_1\)-cycles have equal length.  The possible lengths for
\(\sigma_0\) are \(n\) and \(N\); the length for \(\sigma_1\) is \(m\).
The equality \(n=m\) contradicts coprimality, while \(N=m\) would imply
\(m\mid n\).  Hence no second orbit exists.
\end{proof}

\section{The first quotient Hurwitz problems}
\label{sec:queue}

The complete-chain algorithm of
\cite{guccioneEtAl2017algorithms,guccioneEtAl2022degree108} supplies a
finite queue of terminal data at each degree bound.  We apply
\cref{thm:lattice-gap} to the first five rows at and beyond maximum degree
\(125\).

\begin{table}[ht]
\centering
\small
\begin{tabular}{@{}llcrrlr@{}}
\toprule
case & terminal \(A;k\) & \((m,n)\) & \(g\) & quotient
& quotient passport & classes\\
\midrule
\(F_2\), 125
& \((7/5,2);1\) & \((3,5)\) & 5 & 6
& \((5,1),(3^2),(3,1^3)\) & 1\\
one-step, 126
& \((19/7,5);2\) & \((2,3)\) & 3 & 10
& \((3^3,1),(2^5),(8,1^2)\) & 1\\
two-step, 126
& \((11/6,3);1\) & \((3,2)\) & 2 & 9
& \((2^4,1),(3^3),(7,1^2)\) & 1\\
\(F_{24}\), 128
& \((19/8,3);1\) & \((3,4)\) & 4 & 9
& \((4^2,1),(3^3),(5,1^4)\) & 2\\
one-step, 132
& \((19/4,8);1\) & \((2,3)\) & 3 & 16
& \((3^5,1),(2^8),(13,1^3)\) & 2\\
\bottomrule
\end{tabular}
\caption{The first lattice-compatible terminal Hurwitz problems.}
\label{tab:queue}
\end{table}

\begin{proposition}[Exact Hurwitz count]
\label{prop:hurwitz-count}
The final column of \cref{tab:queue} is exact.  In every row the deck group
of every connected cover is trivial, so the weighted Hurwitz count equals
the ordinary number of isomorphism classes.
\end{proposition}

\begin{proof}
For conjugacy classes \(C_0,C_1,C_\infty\subseteq S_D\), the Frobenius
character formula gives the weighted number
\[
\frac{|C_0||C_1||C_\infty|}{(D!)^2}
\sum_{\chi\in\operatorname{Irr}(S_D)}
\frac{\chi(C_0)\chi(C_1)\chi(C_\infty)}{\chi(1)}.
\]
The exact verifier evaluates the irreducible characters by the
Murnaghan--Nakayama rule.  The five resulting rational numbers are
\(1,1,1,2,2\).

By \cref{thm:lattice-gap}, every triple is transitive.  In all five rows
\(N=1\), so the first branch permutation has a unique fixed sheet.  A deck
transformation centralizing the transitive monodromy group must preserve
that sheet and is therefore the identity.  Thus there is no automorphism
weight to remove.
\end{proof}

\subsection{The degree-\texorpdfstring{\(125\)}{125} quotient}

For the first row, the ambient passport
\[
(5^6),\qquad(3^{10}),\qquad(15,1^{15})
\]
has eleven connected dessin classes: eight with trivial deck group, two
with deck group \(C_3\), and one with deck group \(C_5\).  The gap is five,
so only the last class descends to the polynomial lattice.

\begin{proposition}
\label{prop:f2-map}
Up to source and target normalization, the unique lattice-compatible
quotient is obtained from
\[
\bar p=1-u,\qquad
\bar q=\frac15-\frac35u+\frac9{25}u^2.
\]
It satisfies
\[
\bar p\bar q-3u\bar p\bar q'+5u\bar p'\bar q=\frac15,
\]
and its Belyi map is
\begin{equation}
\label{eq:f2-map}
\bar\tau(u)\doteq
\frac{u(u-1)^5}
{\left(u^2-\frac53u+\frac59\right)^3}.
\end{equation}
The ambient face is the cyclic pullback \(\bar\tau(z^5)\).
\end{proposition}

\begin{proof}
Substitution proves the displayed differential equation.  The derivative
identity and passport follow from \cref{thm:lattice-gap}.  The exact
character and quotient-orbit enumeration proves uniqueness.
\end{proof}

\subsection{Explicit post-\texorpdfstring{\(125\)}{125} covers}

For the one-step degree-\(126\) row, put
\[
P=u^3+u^2+\frac5{12}u+\frac1{18},
\]
\[
Q=u^5+\frac32u^4+u^3+\frac13u^2+\frac5{96}u+\frac1{576}.
\]
The exact identity
\begin{equation}
\label{eq:126-one}
uP^3-Q^2=-\frac{36u^2+28u+9}{2985984}
\end{equation}
gives \(\bar p=18P\), \(\bar q=192Q\).

For the two-step degree-\(126\) row, one may take
\[
\bar p=1+\frac{20}{3}u+24u^2+\frac{288}{7}u^3+\frac{288}{7}u^4,
\]
\[
\bar q=\frac12+5u+12u^2+18u^3.
\]
Equivalently,
\[
P=u^4+u^3+\frac7{12}u^2+\frac{35}{216}u+\frac7{288},
\qquad
Q=u^3+\frac23u^2+\frac5{18}u+\frac1{36}
\]
satisfy
\begin{equation}
\label{eq:126-two}
uP^2-Q^3=-\frac{72u^2+39u+16}{746496}.
\end{equation}

For \(F_{24}\) there are two conjugate solutions over
\(\mathbb Q(\sqrt6)\).  With \(\varepsilon=\pm1\),
\[
\bar p_\varepsilon
=1+u+\left(\frac13+\varepsilon\frac{\sqrt6}{18}\right)u^2,
\]
\[
\bar q_\varepsilon
=\frac14+\frac58u
+\left(\frac25+\varepsilon\frac{\sqrt6}{40}\right)u^2
+\left(\frac{17}{160}
+\varepsilon\frac{11\sqrt6}{480}\right)u^3.
\]
Direct substitution verifies \eqref{eq:quotient-ode} in each case.

\subsection{Infinitesimal rigidity}

Fix the constant terms of a quotient solution and consider
\[
 \mathcal F(p,q)=Npq-mupq'+nup'q.
\]
For \(A=\deg p\) and \(B=\deg q\), its differential is the map
\begin{equation}
\label{eq:face-linearization}
\begin{aligned}
 \mathscr L:
 &\;z k[z]_{\le A}\oplus z k[z]_{\le B}
 \longrightarrow z k[z]_{\le A+B-1},\\
 \mathscr L(\alpha,\beta)
 &=N(\alpha q+p\beta)
 -mu(\alpha q'+p\beta')
 +nu(\alpha'q+p'\beta).
\end{aligned}
\end{equation}
The top coefficient vanishes because \(N-mB+nA=0\), and the constant
coefficient vanishes by the chosen normalization.  Source scaling gives the
kernel vector \((\alpha,\beta)=(up',uq')\).

\begin{proposition}[Rigidity of the explicit reduced covers]
\label{prop:face-rigidity}
For the degree-\(6\), degree-\(10\), degree-\(9\), and two conjugate
degree-\(9\) maps above, \(\mathscr L\) is surjective onto the stated
codomain and its kernel is exactly the source-scaling line.  Hence the
coefficient solution scheme is smooth of dimension one at each point; its
intersection with any affine normalization transverse to source scaling is
a reduced isolated point.
\end{proposition}

\begin{proof}
The exact coefficient matrices have the following ranks and certified
nonzero maximal minors.
\[
\begin{array}{@{}lccc@{}}
\toprule
\text{map} & (\deg p,\deg q) & \operatorname{rank}\mathscr L
& \text{nonzero maximal minor}\\
\midrule
F_2,\ 125 &(1,2)&2&-36/5\\
\text{one-step }126 &(3,5)&7&2090188800\\
\text{two-step }126 &(4,3)&6&37791360/7\\
F_{24},- &(2,3)&4&99/20-153\sqrt6/40\\
F_{24},+ &(2,3)&4&99/20+153\sqrt6/40\\
\bottomrule
\end{array}
\]
In every row the domain dimension is one more than the displayed rank.
The vector \((up',uq')\) is nonzero and lies in the kernel by
differentiating source scaling.  It therefore spans the kernel.
\end{proof}

\begin{remark}
\Cref{prop:face-rigidity} is a coefficient-level statement for the reduced
cover.  It does not prove that a full boundary deformation functor splits
étale-locally into this Hurwitz point and an independent normal-jet problem.
\end{remark}

\section{The normal-boundary determinant operator}
\label{sec:normal-complex}

Let \(E\) be a smooth boundary curve with local coordinate \(z\), and let
\(s\) be a normal parameter.  For positive integers \(\alpha,\beta\),
consider
\begin{equation}
\label{eq:master}
\alpha AB_z-\beta A_zB
+s(A_zB_s-A_sB_z)=\Psi(z),
\end{equation}
where
\[
A=A_0+\sum_{r\ge1}s^ra_r,\qquad
B=B_0+\sum_{r\ge1}s^rb_r.
\]
The order-zero equation is
\[
\alpha A_0B_0'-\beta A_0'B_0=\Psi.
\]

\begin{proposition}[Universal layer operator]
\label{prop:layer}
At normal order \(r\), the terms involving the new pair
\((a_r,b_r)\) are
\begin{equation}
\label{eq:layer}
\boxed{
\mathscr D_r^{\alpha,\beta}(a,b)
=(\alpha-r)a\,dB_0-\beta B_0\,da
+\alpha A_0\,db+(r-\beta)b\,dA_0.}
\end{equation}
Every other term at that order depends only on lower layers.
\end{proposition}

\begin{proof}
The coefficient of \(s^r\) in the first two terms of
\eqref{eq:master} contributes
\[
\alpha(aB_0'+A_0b')-\beta(a'B_0+A_0'b).
\]
The new part of the last term is \(rA_0'b-raB_0'\).  Combining these
terms and multiplying by \(dz\) gives \eqref{eq:layer}.
\end{proof}

\begin{proposition}[Residue adjoint]
\label{prop:adjoint}
For a meromorphic scalar \(\lambda\),
\begin{equation}
\label{eq:adjoint}
\boxed{
(\mathscr D_r^{\alpha,\beta})^\vee(\lambda)
=\left(
\beta B_0\,d\lambda+(\alpha+\beta-r)\lambda\,dB_0,
-\alpha A_0\,d\lambda+(r-\alpha-\beta)\lambda\,dA_0
\right).}
\end{equation}
More precisely,
\begin{align}
\lambda\mathscr D_r(a,b)
&-a\bigl((\mathscr D_r)^\vee\lambda\bigr)_1
-b\bigl((\mathscr D_r)^\vee\lambda\bigr)_2\notag\\
&=d\bigl(-\beta\lambda B_0a+\alpha\lambda A_0b\bigr).
\label{eq:ibp}
\end{align}
Consequently every principal part whose adjoint components annihilate the
allowed correction windows gives a local residue functional annihilating the
image of \(\mathscr D_r\).
\end{proposition}

\begin{proof}
Expand the right side of \eqref{eq:ibp}.  The residue of an exact
meromorphic differential vanishes at every point.
\end{proof}

The precise finite-window statement requires projecting the formal adjoint
to the duals of the allowed correction spaces.  It is proved, with an
explicit coefficient-to-principal-part dictionary, in
\cref{sec:filtered-principal-parts}.

\begin{proposition}[Virtual obstruction excess on \(\PP^1\)]
\label{prop:index}
Suppose
\[
\mathscr D_r:
H^0(\PP^1,\mathcal O(d_A))\oplus
H^0(\PP^1,\mathcal O(d_B))
\longrightarrow
H^0(\PP^1,\mathcal O(d_W)),
\]
with \(d_A,d_B,d_W\ge-1\).  The target-minus-domain dimension is
\begin{equation}
\label{eq:index}
\epsilon=d_W-d_A-d_B-1.
\end{equation}
If \(\mathscr D_r\) is injective, its cokernel has dimension \(\epsilon\).
\end{proposition}

\begin{proof}
Use \(h^0(\mathcal O(d))=d+1\) for \(d\ge-1\).
\end{proof}

For the independently certified full \((8,28)\) support, the windows at
\(5\le r\le8\) are
\[
d_A=10-r,\qquad d_B=15-r,\qquad d_W=25-r.
\]
Thus, for \(5\le r\le8\),
\[
\epsilon_r=r-1.
\]
The exact layer matrices have full column rank in this range, so the virtual
excess is the actual cokernel dimension.

\section{Gauge normal form and filtered principal-part obstructions}
\label{sec:filtered-principal-parts}

Retain the notation
\[
 A=A_0,\qquad B=B_0,\qquad
 \Omega=2A\,dB-3B\,dA=z^2\,dz.
\]
For a layer variation $(a,b)$ put
\[
 u=\frac aA,\qquad v=\frac bB,\qquad
 \xi=u+v,\qquad \eta=2v-3u.
\]

\begin{proposition}[Gauge normal form]
\label{prop:gauge-normal-form}
For every normal order $r$,
\begin{equation}
\label{eq:gauge-normal-form}
 \boxed{
 \mathscr D_r(a,b)
 =AB\left(d\eta+\frac r5\eta\,d\log(AB)\right)
 +\frac{5-r}{5}\xi\,\Omega.}
\end{equation}
Consequently, over the rational function field of the boundary curve,
$\mathscr D_r$ is surjective for $r\ne5$.  At the resonant order,
\begin{equation}
\label{eq:resonant-exact}
 \boxed{\mathscr D_5(a,b)=d(2Ab-3Ba).}
\end{equation}
\end{proposition}

\begin{proof}
Substitute
\[
 u=\frac{2\xi-\eta}{5},\qquad
 v=\frac{3\xi+\eta}{5}
\]
into
\[
 \mathscr D_r(a,b)
 =(2-r)a\,dB-3B\,da+2A\,db+(r-3)b\,dA
\]
to obtain \eqref{eq:gauge-normal-form}.  If $r\ne5$, set $\eta=0$ and
$\xi=5\Phi/((5-r)\Omega)$ to solve $\mathscr D_r(a,b)=\Phi$ over the
function field.  For $r=5$, the $\xi$ term disappears and
$AB(d\eta+\eta\,d\log(AB))=d(AB\eta)=d(2Ab-3Ba)$.
\end{proof}

\begin{remark}[An extended deformation complex]
Allowing a variation $\rho\Omega$ of the normalized right-hand side, put
\[
 \widetilde{\mathscr D}_r(a,b,\rho)
 =\mathscr D_r(a,b)-\rho\Omega.
\]
Then
\[
 G_r(h)=(2Ah,3Bh,(5-r)h)
\]
satisfies $\widetilde{\mathscr D}_r\circ G_r=0$.  Thus the fixed layer is
represented by the three-term complex
\[
 \mathcal G_r\xrightarrow{G_r}
 \mathcal U_{A,r}\oplus\mathcal U_{B,r}\oplus\mathcal R_r
 \xrightarrow{\widetilde{\mathscr D}_r}\mathcal W_r.
\]
At $r=5$ the right-hand-side component of $G_r$ vanishes, leaving the
residual normal-coordinate gauge in the fixed-volume problem.
\end{remark}

\subsection{The exact two-point windows}

Let $o=(z=0)$ and $p=(z=\infty)$ on $E\simeq\mathbf P^1$.  For
$5\le r\le8$, the full Newton support gives
\begin{align}
 U_{A,r}&=\langle1,z,\ldots,z^{10-r}\rangle
        =H^0(E,\mathcal O((10-r)p)),\label{eq:UA-window}\\
 U_{B,r}&=\langle1,z,\ldots,z^{15-r}\rangle
        =H^0(E,\mathcal O((15-r)p)),\label{eq:UB-window}\\
 W_r&=\langle z^{-1},1,z,\ldots,z^{24-r}\rangle\,dz
     =H^0(E,\Omega_E^1(o+(26-r)p)).\label{eq:W-window}
\end{align}
The exact layer matrices have full column rank, hence
\[
 \dim\operatorname{coker}\mathscr D_r
 =(26-r)-(11-r)-(16-r)=r-1.
\]

\begin{proposition}[Filtered principal-part duality]
\label{prop:filtered-principal-part-duality}
Order the target basis as $z^I dz$, $-1\le I\le24-r$.  If
$c=(c_I)$ is a row vector, set
\[
 \lambda_c(z)=\sum_{I=-1}^{24-r}c_Iz^{-I-1}.
\]
Then $c$ lies in the left kernel of the layer matrix if and only if
\begin{align}
 \operatorname{Res}_{0}\!
 \left(z^i\bigl(3B\,d\lambda_c+(5-r)\lambda_c\,dB\bigr)\right)&=0,
 &&0\le i\le10-r,\label{eq:filtered-adjoint-A}\\
 \operatorname{Res}_{0}\!
 \left(z^j\bigl(-2A\,d\lambda_c+(r-5)\lambda_c\,dA\bigr)\right)&=0,
 &&0\le j\le15-r.\label{eq:filtered-adjoint-B}
\end{align}
For the lower-layer forcing form $\Phi_r$, the corresponding compatibility
coordinate is
\begin{equation}
\label{eq:residue-coordinate}
 c^T(-\Phi_r)
 =\operatorname{Res}_{0}\bigl(\lambda_c(-\Phi_r)\bigr).
\end{equation}
\end{proposition}

\begin{proof}
The identity
\begin{align*}
 \lambda\mathscr D_r(a,b)
 &-a\bigl(3B\,d\lambda+(5-r)\lambda\,dB\bigr)\\
 &-b\bigl(-2A\,d\lambda+(r-5)\lambda\,dA\bigr)
 =d(-3\lambda Ba+2\lambda Ab).
\end{align*}
shows that the residue of the left side is zero.  Taking successively
$a=z^i,b=0$ and $a=0,b=z^j$ gives
\eqref{eq:filtered-adjoint-A}--\eqref{eq:filtered-adjoint-B}.  Conversely,
those equalities say exactly that the residue functional annihilates every
column of the ordered coefficient matrix.  Formula
\eqref{eq:residue-coordinate} is the definition of the dual monomial basis.
\end{proof}

\begin{remark}
Away from resonance the two adjoint expressions in
\eqref{eq:filtered-adjoint-A}--\eqref{eq:filtered-adjoint-B} need not vanish
as meromorphic differentials.  Only their projections to the duals of the
allowed Newton windows vanish.  The obstruction classes are therefore
filtered principal parts, rather than unrestricted global adjoint
solutions.
\end{remark}

\subsection{The exact cokernel frames through order eight}

The deterministic row reduction gives echelon principal-part frames whose
leading pole orders at $o$ are
\[
\begin{array}{c|c|c}
 r&\dim\operatorname{coker}\mathscr D_r&
 \text{leading pole orders}\\
\hline
5&4&0,18,19,20\\
6&5&0,16,17,18,19\\
7&6&0,14,15,16,17,18\\
8&7&0,12,13,14,15,16,17.
\end{array}
\]
In every layer the order-zero class is represented by $\lambda=1$ and its
forcing pairing vanishes identically.  At $r=5$, the order-twenty class also
pairs identically to zero; the order-nineteen pairing vanishes after the
order-four normalization.  At $r=6$, the order-nineteen pairing vanishes
after that normalization.  The surviving normalized equations are therefore
\[
1,3,5,6
\]
at orders $5,6,7,8$, respectively.

In particular, the five weight-seven equations are exactly the residue
coordinates with leading pole orders $14,15,16,17,18$, and the six
weight-eight equations are exactly those with leading pole orders
$12,13,14,15,16,17$.  The exact principal parts and normalization scalars
are contained in the accompanying machine-readable certificate.

\begin{remark}[The six resultant coordinates]
The sparse resultant and norm certificate uses
\[
 (F_4,F_6,F_8,F_9,F_{10},F_{11}).
\]
In the filtered frames above these are the layer-seven coordinates of
leading pole orders $14,16,18$ and the layer-eight coordinates of leading
pole orders $12,13,14$.  This is an exact coordinate projection of the full
fifteen-component obstruction section.  It is not a whole layer or a
subspace defined by one pole cutoff.  The present calculation therefore
does not identify this six-dimensional projection as a canonical
subquotient of the pole filtration; its current justification is its exact
unit-ideal and mixed-volume effectiveness.
\end{remark}


\section{A universal secondary Belyi map}
\label{sec:secondary}

Let \(D=\{t=0\}\) be a boundary component with coordinate \(s\), and put
\[
c(s)=\frac{s}{s-1}.
\]
Assume target coordinates have been normalized to
\begin{equation}
\label{eq:contact}
\pi=t^ac(s)+O(t^{a+1}),\qquad
u=\sum_{j\ge1}t^jh_j(s),
\end{equation}
where \(a\ge1\).  Suppose the boundary-volume identity has no term below
order \(t^{a+e-1}\) and that its first nonzero coefficient is
\begin{equation}
\label{eq:volume}
d\pi\wedge du
=-t^{a+e-1}(s-1)^{-a-2}\,ds\wedge dt
+O(t^{a+e}).
\end{equation}

At a lower order \(j\), the homogeneous equation is
\begin{equation}
\label{eq:homogeneous}
jc'h_j-ach_j'=0.
\end{equation}

\begin{lemma}[Resonance and target shears]
\label{lem:resonance}
A nonzero rational solution of \eqref{eq:homogeneous} exists if and only if
\(a\mid j\).  If \(j=ak\), every solution is \(h_j=Cc^k\), and its
leading contribution \(Ct^{ak}c^k\) is removed by the target shear
\(u\mapsto u-C\pi^k\).
\end{lemma}

\begin{proof}
Equation \eqref{eq:homogeneous} gives \(h_j=Cc^{j/a}\).  Since
\(\operatorname{div}(c)=(0)-(1)\), the power is rational exactly when
\(j/a\) is an integer.  The assertion about the shear follows from
\eqref{eq:contact}.
\end{proof}

After removing all resonant lower terms, the coefficient of
\eqref{eq:volume} gives
\begin{equation}
\label{eq:transport-equation}
ec'h-ach'=-(s-1)^{-a-2}.
\end{equation}

\subsection{The contact degree}

\begin{proposition}[Contact-degree formula]
\label{prop:contact}
Suppose coprime integers \(m,n\) satisfy
\[
v_E(P)=-m,\quad v_E(Q)=-n,\qquad
v_D(P)=-ma,\quad v_D(Q)=-na.
\]
Choose \(c,d\in\mathbb Z\) with \(dn-cm=1\) and set
\[
\pi=P^c/Q^d,\qquad \tau=Q^m/P^n.
\]
Assume that in the toric chart
\[
 x=t^{-1}\varepsilon_x(t,s),\qquad
 dx\wedge dy=t^{-2}\varepsilon_\omega(t,s)\,dt\wedge ds,
\]
where both \(\varepsilon_x\) and \(\varepsilon_\omega\) are units, and that
\[
 dP\wedge dQ=x^\kappa\varepsilon(s,t)\,dx\wedge dy
\]
with \(\varepsilon\) a unit.  After removing all lower resonant shears,
assume
\(\tau-\tau|_D=t^e h(s)+O(t^{e+1})\) with nonzero leading wedge
coefficient.  Then the first possible non-shear contact degree is
\begin{equation}
\label{eq:contact-degree}
e=a(m+n)-\kappa-1.
\end{equation}
\end{proposition}

\begin{proof}
The Bézout relation gives \(v_D(\pi)=a\) and \(v_D(\tau)=0\).  Direct
differentiation gives
\[
d\pi\wedge d\tau
=-P^{c-n-1}Q^{m-d-1}\,dP\wedge dQ.
\]
The monomial factor has \(t\)-order \(a(m+n+1)\), while the last wedge has
order \(-\kappa-2\).  On the other hand, the wedge of
\(\pi\sim t^a\) with \(\tau-\tau|_D\sim t^e\) has order \(a+e-1\).
Equating orders yields \eqref{eq:contact-degree}.
\end{proof}

\subsection{Transport and passport}

\begin{theorem}[Universal secondary-boundary transport]
\label{thm:secondary}
Let \(a,e\in\mathbb N\) satisfy \(e>a^2\), and put
\(\delta=e-a^2\).  Equation \eqref{eq:transport-equation} has a unique
solution
\[
h(s)=\frac{H_{a,e}(s)}{(s-1)^a}
\]
whose numerator is a polynomial of degree exactly \(a\), where
\begin{align}
H_{a,e}(s)
&=\sum_{k=0}^a
\frac{a^k a!}{(a-k)!\prod_{j=0}^k(e-aj)}s^k
\label{eq:H-sum}\\
&=\frac1e\,{}_2F_1\!\left(-a,1;1-\frac ea;s\right).
\label{eq:H-hypergeometric}
\end{align}
It satisfies
\begin{equation}
\label{eq:H-ode}
as(s-1)H'+(e-a^2s)H=1,
\end{equation}
\[
H(0)=\frac1e,\qquad H(1)=\frac1\delta,
\]
and \(H\) is squarefree.

The exceptional ratio
\begin{equation}
\label{eq:secondary-W}
W_{a,e}(s)=\frac{(s-1)^\delta H_{a,e}(s)^a}{s^e}
\end{equation}
is a degree-\(e\) Belyi map.  It obeys
\begin{equation}
\label{eq:secondary-derivative}
W_{a,e}'(s)=
\frac{(s-1)^{\delta-1}H_{a,e}(s)^{a-1}}{s^{e+1}},
\end{equation}
and its passport is
\begin{equation}
\label{eq:secondary-passport}
(a^a,\delta),\qquad(e),\qquad(a+1,1^{e-a-1}).
\end{equation}
\end{theorem}

\begin{proof}
Substitute \(h=H/(s-1)^a\) into
\eqref{eq:transport-equation}; since \(c'=-(s-1)^{-2}\), this gives
\eqref{eq:H-ode}.  Writing \(H=\sum h_ks^k\) gives
\[
h_0=\frac1e,\qquad
h_k=\frac{a(a+1-k)}{e-ak}h_{k-1}.
\]
The denominators are nonzero because \(e>a^2\), and the recurrence gives
\eqref{eq:H-sum}.  It also proves uniqueness.  The hypergeometric notation
is the same terminating recurrence.  Evaluation at zero and one gives the
stated endpoint values.  If \(H(s_0)=0\), then
\eqref{eq:H-ode} gives \(as_0(s_0-1)H'(s_0)=1\); hence every root is simple
and avoids \(0,1\).

Logarithmic differentiation of \eqref{eq:secondary-W}, followed by
\eqref{eq:H-ode}, gives
\[
\frac{W'}W=\frac1{s(s-1)H},
\]
which is \eqref{eq:secondary-derivative}.  The numerator and denominator of
\(W\) are coprime of degree \(e\).  Over zero, the \(a\) roots of \(H\)
have multiplicity \(a\) and \(s=1\) has multiplicity \(\delta\).  Over
infinity, \(s=0\) has multiplicity \(e\).  At \(s=\infty\),
\eqref{eq:secondary-derivative} shows that
\(W-W(\infty)\) has order \(a+1\).  Finally,
\[
a(a-1)+(\delta-1)+(e-1)+a=2e-2,
\]
so Riemann--Hurwitz leaves only \(e-a-1\) unramified points in the third
fiber and proves \eqref{eq:secondary-passport}.
\end{proof}


\begin{remark}[Truncated-binomial Shabat form]
Let \(\delta=e-a^2\), let \(h_a=[s^a]H_{a,e}(s)\), and put
\[
 T_{a,\delta}(x)=\frac{x^aH_{a,e}(1/x)}{h_a}.
\]
Reversing the coefficient recurrence gives
\[
 T_{a,\delta}(x)=
 \sum_{j=0}^a\frac{(\delta/a)_j}{j!}x^j,
\]
and
\[
 \frac{W_{a,e}(1/x)}{W_{a,e}(\infty)}
 =(1-x)^\delta T_{a,\delta}(x)^a.
\]
Thus \(W_{a,e}\) is the reciprocal form of the classical
truncated-binomial Shabat family; compare \cite{vidunas2025belyi}.  The
content of \cref{thm:secondary} for the present application is that the
Keller boundary-volume equation forces this family and determines
\((a,e)\) from the adjacent valuation data.
\end{remark}

\begin{corollary}[The \((4,17)\) cover]
\label{cor:417}
For \(a=4,e=17\),
\[
H_{4,17}(s)
=\frac{195+240s+320s^2+512s^3+2048s^4}{3315},
\]
\[
W_{4,17}(s)=\frac{(s-1)H_{4,17}(s)^4}{s^{17}},
\qquad
W_{4,17}'(s)=\frac{H_{4,17}(s)^3}{s^{18}},
\]
with passport \((4^4,1),(17),(5,1^{12})\).
\end{corollary}

\section{Filtered descent and the global gluing problem}
\label{sec:descent}

The phrase ``vanishing high-pole residues lowers the Newton window'' has an
exact linear core.

\begin{lemma}[Filtered one-layer descent]
\label{lem:filtered}
Let \(D:U\to V\) be a linear map of finite-dimensional vector spaces and
let \(V_{\le d}\subseteq V\) be a filtered subspace.  For
\(\Phi\in V\), the following are equivalent:
\begin{enumerate}[label=(\roman*)]
\item there exists \(u\in U\) such that
\(\Phi-Du\in V_{\le d}\);
\item every
\[
\lambda\in(\operatorname{im}D)^\perp
\cap(V_{\le d})^\perp
\]
satisfies \(\lambda(\Phi)=0\).
\end{enumerate}
\end{lemma}

\begin{proof}
Let \(q:V\to V/\operatorname{im}D\).  Condition (i) says
\(q(\Phi)\in q(V_{\le d})\).  Finite-dimensional duality says precisely
that every functional on the quotient vanishing on \(q(V_{\le d})\) must
also vanish on \(q(\Phi)\).  Pulling the functionals back to \(V\) gives
(ii).
\end{proof}

For \(\mathscr D_r\), the relevant functionals are the pole-filtered
adjoint residue classes of \cref{prop:adjoint}.  Thus at one fixed normal
layer their vanishing is equivalent to reducing the forcing term to the
smaller pole window.

This is not yet a descent theorem for polynomial Keller pairs.  The
corrections must be chosen compatibly at all layers, and the resulting
formal transformation must integrate to an allowable polynomial
approximate-root operation.

\begin{question}[Integrable filtered descent]
\label{q:integrable}
For a terminal complete-chain boundary model, does simultaneous vanishing
of all high-pole residue obstructions either force the boundary-gluing
scheme to be empty or produce a polynomial transformation of strictly
smaller complete-chain complexity?
\end{question}

A positive answer, together with completeness of the finite boundary model,
would permit induction on complete-chain complexity.  Neither hypothesis is
proved here.

\begin{question}[Normal splitting]
\label{q:splitting}
Near one of the rigid reduced covers of
\cref{prop:face-rigidity}, does the complete boundary deformation functor
split étale-locally into the reduced Hurwitz point and a normal-jet
deformation problem?  If not, which cross-terms measure the failure?
\end{question}

\input{sections/toric-norm}

\section{Exact computation and reproducibility}
\label{sec:computation}

All finite computations in this paper use exact arithmetic.  The archived
certificate suite performs the following checks.

\begin{enumerate}[label=(\arabic*)]
\item It verifies the universal layer operator, the integration-by-parts
identity, the logarithmic gauge normal form, and
\(\dim\operatorname{coker}\mathscr D_r=r-1\) for \(5\le r\le8\).
\item It reconstructs every layer matrix, forcing vector, left-null basis,
and principal-part representative through order eight; verifies the
filtered adjoint identities; and reproduces the fifteen normalized
compatibility equations coefficient-for-coefficient.
\item For \(1\le a\le8\) and \(\delta\in\{1,2,3\}\), it verifies the
secondary recurrence, squarefreeness, derivative identity, degree,
passport, and Riemann--Hurwitz count for \(e=a^2+\delta\).
\item It verifies the terminal-corner arithmetic, lattice divisibilities,
quotient equations, derivative formulas, and explicit quotient maps.
\item It evaluates the exact symmetric-group character sums, recovers the
eleven ambient degree-\(30\) classes, selects the unique lattice-compatible
\(C_5\) pullback, and verifies the five quotient counts \(1,1,1,2,2\).
\item It constructs the exact tangent matrices for the explicit quotient
points, proves their ranks, checks their scaling kernels, and records a
nonzero maximal minor in each case.
\item For the terminal six-coordinate subsystem, it verifies the exact
mixed-volume identity
\[
 \operatorname{MV}(A,B,C,C,C)=296.
\]
It enumerates all \(344\) proper toric faces, proves the \(74\) nontrivial
saturated initial ideals are unit ideals over \(\mathbf F_{2053}\), and
verifies the complete \(296\)-dimensional finite \'{e}tale algebra.
\item It checks the coordinate units, the squarefree primitive-element
polynomial, \(\det(m_\rho)=682\) at the chosen split prime, and the five
Galois-conjugate determinant residues whose rational norm is \(51\) modulo
\(2053\).
\end{enumerate}

Freshly generated machine-readable outputs agree with the archived
references.  These computations prove the exact face, residue, Hurwitz, and
degree-\(296\) terminal-model statements made in this paper.  They do not
verify that the imported complete-chain reduction, support enumeration,
normalizations, and saturations exhaust every hypothetical counterexample,
and they do not establish novelty.

\appendix
\input{appendices/descent-puiseux-log}
\section{The degree-\(21\) divisor and terminal certificates}
\label{app:degree-twenty-one}

The body develops the lattice-gap quotient covers.  Here we record the
complementary calculation on the original lower face and the exact
certificates obtained after substituting its five dessins into the two
surviving Newton supports.

\subsection{The lower face}

For either surviving support in the Guccione--Guccione--Horruitiner--Valqui
reduction \cite{guccioneEtAl2022degree108}, the transformed bracket is
\[
[P,Q]=X^2.
\]
For the primitive toric valuation
\[
\nu(X)=-2,\qquad \nu(Y)=1,
\]
put \(z=XY^2\).  The lower faces have the form
\[
P_{\mathrm{face}}=Xp(z),\qquad
Q_{\mathrm{face}}=X^2Yq(z),
\]
where \(\deg p=7\), \(\deg q=10\), and
\begin{equation}
\label{eq:degree21-face}
pq+2zpq'-3zp'q=1.
\end{equation}

\begin{proposition}[The degree-\(21\) Belyi map]
\label{prop:degree21-belyi}
The rational function
\[
\tau(z)=z\frac{q(z)^2}{p(z)^3}
\]
is a degree-\(21\) Belyi map with passport
\[
(2^{10}1),\qquad(3^7),\qquad(17\,1^4).
\]
There are exactly five connected dessin isomorphism classes with this
passport.  They form one arithmetic orbit over an irreducible quintic field
and have monodromy \(A_{21}\).
\end{proposition}

\begin{proof}
Equation \eqref{eq:degree21-face} implies that \(p,q\) are coprime and
nonzero at zero.  Differentiation gives
\[
\tau'(z)=\frac{q(z)}{p(z)^4}.
\]
Thus the ten roots of \(q\) give double zeros, the seven roots of \(p\)
give triple poles, and infinity has ramification index \(17\).  The total
ramification is
\[
10+14+16=40=2\cdot21-2,
\]
so there is no additional branch value.

The exact Murnaghan--Nakayama character calculation in \(S_{21}\), followed
by the class-size factor, gives \(5\cdot21!\) labeled triples.  Each is
transitive, and the deck group is trivial, so there are five isomorphism
classes.  Exact coefficient reconstruction places all five over one
quintic field; exact permutation certificates give monodromy \(A_{21}\).
The enumeration and reconstruction are computer-assisted.
\end{proof}

\subsection{Identification with Borisov's divisor}

\begin{theorem}[Valuation over \(F_{-5}\)]
\label{thm:borisov-fminusfive}
Let \(E_\nu\) be the source divisor defined by \(\nu\).  Its target center is
Borisov's divisor \(F_{-5}\).  The normal ramification index and residue
degree are
\[
(e,f)=(1,21).
\]
Consequently neither surviving Newton support can realize Borisov's
Three-dessin framework.
\end{theorem}

\begin{proof}
Set \(t=Y\), so \(x=t^2/z\) in the original source coordinates.  The source
volume form is
\[
dx\wedge dy=t^{-6}z^2\,dt\wedge dz.
\]
After adding the reduced boundary coefficient, the augmented-canonical label
is \(-5\).

The face expansions are
\[
P=t^{-2}zp(z)+\text{higher \(\nu\)-terms},\qquad
Q=t^{-3}z^2q(z)+\text{higher \(\nu\)-terms}.
\]
On the target put
\[
s=\frac P Q,\qquad \lambda=\frac{Q^2}{P^3}.
\]
Then \(s\) is a uniformizer and
\[
dP\wedge dQ=-s^{-6}\lambda^{-3}\,ds\wedge d\lambda,
\]
so the target label is also \(-5\).  Pullback gives
\[
\frac P Q
=t\left(\frac{p(z)}{zq(z)}+O(t)\right),
\]
hence \(e=1\).  The residue of \(\lambda\) is
\[
\operatorname{res}_{E_\nu}(\lambda)
=z\frac{q(z)^2}{p(z)^3}=\tau(z),
\]
so \(f=21\) by \cref{prop:degree21-belyi}.

In Borisov's target coordinates \((y_1,y_2)=(Q,P)\), the divisor
\(F_{-5}\) is characterized by valuations \((-3,-2)\) and by the
nonconstant residual parameter \(y_1^2/y_2^3\).  Those are exactly the
values above.  Later infinitely-near divisors with the same two numerical
orders have constant residue of this parameter, so they are distinguished.

Borisov's Three-dessin framework gives residue degree \(16\) on its unique
source divisor above \(F_{-5}\) \cite{borisov2020frameworks}.  The forced
divisor has residue degree
\(21\), a contradiction.  This excludes the named framework, not the two
supports themselves.
\end{proof}

\begin{corollary}
Any plane counterexample of maximum coordinate degree below \(125\) has
generic degree at least \(21\).
\end{corollary}

\begin{proof}
The published below-\(125\) reduction leaves the two supports above, and the
generic-degree formula contains the contribution \(ef=21\) from
\cref{thm:borisov-fminusfive}.
\end{proof}

\begin{question}
Is \(E_\nu\) the only source divisor over the generic point of \(F_{-5}\)?
The calculation \(f=21\) is one residue-degree contribution; cancellation
valuations or infinitely-near divisors could contribute additional sheets.
\end{question}

\subsection{Exact support certificates}

Substituting the five reconstructed lower faces into the truncated Newton
support gives seven compatibility equations over the quintic field.  A
nonzero exact resultant, certified by good reduction, eliminates every
dessin.  Thus the truncated support is empty in characteristic zero.

For the full support, exact layer elimination and normalization reduce the
problem to fifteen equations in five variables over the same quintic field.
Six selected equations already generate the unit ideal.

\begin{theorem}[Terminal unit-ideal certificate]
\label{thm:terminal-unit}
For the stored normalized full-support system, the six selected obstruction
equations generate the unit ideal over the quintic coefficient field.
\end{theorem}

\begin{proof}[Exact computer-assisted certificate]
The Macaulay multiplication map has \(10{,}824\) columns.  The certificate
specifies \(7{,}121\) rows and \(7{,}121\) pivot columns.  Reducing this
exact minor at a good prime and at the selected quintic embedding gives
\[
\det=859\pmod{2053}.
\]
The determinant is therefore nonzero in characteristic zero.  The
certificate package regenerates all fifteen exact equations, rebuilds the
matrix, and replays the fixed pivot minor.
\end{proof}

Combining the truncated resultant and
\cref{thm:terminal-unit} eliminates both encoded support alternatives after
the degree-\(21\) face substitution.  This statement is exact for the stored
systems.  Turning it into a stand-alone global below-\(125\) proof requires
a line-by-line audit that the published Newton reduction, the five faces,
all saturations and normalizations, and the layer systems exhaust every
candidate.

There is also a smaller geometric compression.  Five partial-gluing
equations have mixed volume \(296\); on the certified
\(296\)-dimensional modular finite algebra, the sixth obstruction is a unit,
and its rational norm over the five split embeddings is
\[
51\pmod{2053}.
\]
This is already a compact representative of the large certificate.  To use
it as an independent characteristic-zero proof, one still needs a
publication-level toric-saturation lemma showing that the finite \'{e}tale
special fiber exhausts the complete toric intersection.

\subsection{Credit for the degree-\(125\) conclusion}

MathOverflow user \texttt{ratto3423} publicly announced the conclusion that
every characteristic-zero plane counterexample has maximum coordinate degree
at least \(125\), based on a computer-assisted elimination of the same two
supports \cite{ratto3423degree125}.  We credit that public theorem
announcement to \texttt{ratto3423}.  The public answer does not disclose
enough method to compare its computation with the certificates above.

Accordingly, our claim is not priority for the degree-\(125\) conclusion.
The potential contribution of this project is the explicit degree-\(21\)
boundary description and an independently auditable terminal certificate
architecture.  The current certificate is conditional only at the upstream
exhaustiveness layer described above, not at its final linear-algebra step.


\section*{AI and computational disclosure}

AI systems were used extensively in exploration, symbolic-program
development, proof drafting, source organization, and manuscript editing.
They are not authors.  Human specialist review has not yet been completed.
The evidence for the computer-assisted propositions is the exact certificate
suite described above, not language-model output.

\bibliographystyle{amsplain}
\bibliography{../common/references}

\end{document}
